%% =====================================================================
%%  CHAPTER 3: PROPOSED SOLUTION AND DATA              (budget 6 pages)
%%  chapters/ch3-solution.tex
%%  Input by \pcIfSolution{afterproblem} (end of ch2) or
%%  \pcIfSolution{afteralternatives} (end of ch6); see \SolutionPlacement.
%% =====================================================================
\pcChapter{solution}{6}{הפתרון המוצע והנתונים}
\begin{pcguide}[הפרק החשוב ביותר: מה אתם מציעים, ועם איזה מידע]
זהו הפרק שאומר \textbf{מה הפתרון, מה נבנה ועם אילו נתונים נעבוד} (מה יש, מאיפה הוא מגיע ומה חסר). זה דוח מכין ולא דוח ביצוע, ולכן
מנסחים בלשון תוכנית (“נבנה”, “נבדוק”). הדוח הוא חוזה מול המרצה: מה תעשו, מה תהיה התוצאה, ואיך ייבדק שעמדתם בכך.
סעיפים~\ref{sec:solstack} ו-\ref{sec:solgen} מפרטים את הטכנולוגיה ואת “הכלי שמייצר את הכלי”.
\end{pcguide}

\section{הצהרת הפתרון}\label{sec:solstate}
\begin{pcbox}[pcGood]{הצהרת הפתרון: משפט אחד}
נציע לפתח אב-טיפוס מקומי של מערכת הדרכה המונעת במנוע \textbf{\latin{Grounded RAG}} עבור יחידת ההדרכה והרכש בתעשייה האווירית. המערכת תקלוט מסמכים הנדסיים
לא מסווגים ולקחי מומחים, תאחזר מהם ראיות לפי מקור וגרסה, ותפיק יחידות לימוד, שאלות ותצורות תרחיש \latin{JSON} שמאכלסות תבניות
\threeD{} פרמטריות, והכול באישור מומחה לפני פרסום. בפיילוט בן \latin{14} שבועות נבדוק: דיוק חילוץ ידע \pctto{60}{90} (בסיס מוערך), קיצור זמן
הפקה של \gesym{\pct{50}}, נאמנות עובדתית של \gesym{\pct{90}} והשלמת משימות \threeD{} של \gesym{\pct{80}}.
\end{pcbox}
\textbf{הערך לארגון:} חפיפה שאינה תלויה בזמינות מומחה בודד, יחידות הדרכה עקיבות למקור, ותרחישי תרגול שמתעדכנים כשהנוהל משתנה.
\textbf{מענה לפער הידע:} השילוב בין הפקה מעוגנת, אישור מומחה, תרגול \threeDadj{} ומדידת זמן מול בסיס ידני לא נבדק במשותף בספרות שנסקרה (סעיף~\ref{sec:gap}).

\section{מושג הפתרון והיקפו}\label{sec:solconcept}
איור~\ref{fig:solconcept} מציג את מושג הפתרון: מקורות מאושרים נכנסים למנוע \latin{RAG} מקומי, ממנו יוצאים תוצרי הדרכה שעוברים אישור מומחה לפני
שהם מגיעים ללומד. מחזור החיים כולל שמונה שלבים:
\begin{enumerate}[nosep]
  \item \textbf{קליטה ובקרה:} מנהל הידע בוחר מקורות מותרים ורושם בעלים, גרסה ורמת הרשאה.
  \item \textbf{הכנת קורפוס:} פענוח, פירוק לפי סעיפים תוך שימור עמוד ומקור, ואינדוקס מקומי (סעיף~\ref{sec:solstack}).
  \item \textbf{אחזור:} חיפוש מילולי וסמנטי המסונן לפי הרשאה ולפי גרסה בתוקף.
  \item \textbf{הפקה:} מודל שפה מקומי מקבל רק קטעים מאושרים ומכין יחידה ושאלות עם סימוכין; בהיעדר מקור הוא מחזיר “חוסר ראיות”.
  \item \textbf{אישור מומחה:} פריט שנדחה חוזר לעיבוד ואינו מתפרסם.
  \item \textbf{הפקת תצורת תרחיש:} קובץ \latin{JSON} שמאכלס תבנית \threeD{} פרמטרית (סעיף~\ref{sec:solgen}).
  \item \textbf{הכשרה ובחינה:} הלומד צופה ביחידה, עונה ונכנס לסצנה; בחירות, ציונים וזמני תגובה נשמרים.
  \item \textbf{למידה ממקרים חדשים:} מקרי קצה מריאיון מומחה נכנסים לתור אישור ואינם הופכים אוטומטית לנוהל.
\end{enumerate}
\begin{figure}[!htbp]
\centering
\pcFig{solution_concept}{%
\pcAutoFit{\begin{pctikz}[x=1cm,y=1cm]
  \node[pcnode,text width=3.2cm] (in) at (0,0) {\tn{קלט: מסמכים פומביים מאושרים ולקחי מומחים (גרסה, עמוד, הרשאה)}};
  \node[pcnodeD,text width=3.6cm] (core) at (-5.2,0) {\tn{מנוע \latin{Grounded RAG} מקומי: אחזור, הפקה עם סימוכין ואישור מומחה}};
  \node[pcnode,text width=3.2cm] (out) at (-10.4,0) {\tn{פלט: יחידת לימוד, מבחן ותצורת \latin{JSON} לתבנית \latin{3D}}};
  \node[pcnode,text width=3.6cm,fill=white,dashed] (ai) at (-5.2,-1.9) {\tn{מודל שפה \latin{(LLM)} מקומי, ללא רשת יוצאת}};
  \draw[pcflow] (in) -- (core); \draw[pcflow] (core) -- (out); \draw[pcflow,dashed] (ai) -- (core);
  \node[font=\scriptsize\sffamily,text=pcGrey] at (-5.2,1.3) {\tn{זרימת ההחלטה: מימין לשמאל}};
\end{pctikz}}}
\caption{מושג הפתרון: מקורות מאושרים, מנוע \latin{Grounded RAG} מקומי ותוצרי הדרכה ותרגול \latin{3D}}\label{fig:solconcept}
\end{figure}
\textbf{היקף והנחות.} בתוך ההיקף: מנוע \latin{RAG} מקומי עם אישור מומחה, מבחן ומחוון; שלוש תבניות \threeD{} פרמטריות וסצנה אחת עם שלוש נקודות החלטה, עם מסלול \twoD{} מקביל;
ומדידת זמן, נאמנות ורישום החלטות. מחוץ להיקף: הפעלה אוטומטית של מערכות ייצור, הסמכה ארגונית של עובד, הפצת מידע מסווג, דגם נאמן של ציוד מסווג,
הדמיה מרובת משתתפים ותרגיל ארגוני פיזי. ההנחות: קורפוס \latin{FAA} פומבי מייצג את מבנה המסמכים הארגוניים, מומחה תוכן זמין לחלונות ביקורת קבועים, ומודל מקומי
בגודל \latin{7B}--\latin{8B} מספיק לניסוי ללא תקשורת יוצאת. העקיבות מדרישה לבדיקה נמצאת בטבלה~\ref{tab:eval} ובפרק~\ref{ch:problem}.

\section{מפרט טכנולוגי: מודל מקומי, אחזור היברידי ומסד וקטורי}\label{sec:solstack}
המפרט שלהלן הוא תכן מוצע לפיילוט. בחירת המודל הסופי תיעשה בשבועות \rng{3}{5} לאחר מדידת נאמנות וקצב על הקורפוס הפומבי, ועד אז כל פריט הוא ברירת
מחדל מנומקת (טבלה~\ref{tab:stack}). כל הרכיבים פועלים על שרת מקומי אחד, ללא תקשורת יוצאת (\req{09}).
\begin{table}[!htbp]
\centering
\pcCompactTables
\caption{מפרט הטכנולוגיה המוצע לפיילוט: בחירה, נימוק ומגבלה}\label{tab:stack}
\begin{pctable}{H{0.15},H{0.33},H{0.34},C{0.12}}
  \pchead{שכבה & בחירה מוצעת & נימוק ומגבלה & מקור}
  \pcrow{חומרה & תחנת עבודה מקומית עם כרטיס \latin{GPU} בעל יותר מ-\latin{16} \latin{GB} זיכרון (מומלץ \latin{24} \latin{GB}); חומרה מבודדת מהרשת & חישוב התכן מראה צורך של כ-\mth{9.4} \latin{GB} (\mth{\approx12.6} \latin{GB} בהקשר ארוך); יימדד בשבוע \latin{5} & טבלה~\ref{tab:budget}}
  \pcrow{מודל שפה & מודל בעל משקולות פתוחות בגודל \latin{7B}--\latin{8B}, מכוון להוראות, מכומת ל-\latin{4} סיביות (למשל משפחות \latin{Llama 3.1 8B Instruct} או \latin{Qwen2.5 7B Instruct}; שני מועמדים ייבחנו) & משקלים של כ-\mth{4} \latin{GB}; פועל ללא ענן. השפעת הקוונטיזציה על הנאמנות תימדד ולא תוסק מהספרות; תמיכת עברית בפלט תיבדק כשער נפרד, והקורפוס באנגלית & \cite{lin2024awq}}
  \pcrow{הגשת המודל & \latin{Ollama} בפיתוח; \latin{vLLM} לבדיקת קצב ובקשות מקבילות & ניהול זיכרון מטמון \latin{KV} בשיטת \latin{PagedAttention} מעלה תפוקה; יעד הקצב \gesym{37} טוקנים לשנייה (\req{08}) & \cite{kwon2023vllm}}
  \pcrow{מודל ייצוג (\latin{Embedding}) & מודל רב-לשוני, למשל \latin{BGE-M3} (למעלה מ-\latin{100} שפות, קלט ארוך, ייצוג צפוף ודליל) & מאפשר שאילתות בעברית על מסמכים באנגלית; מועמד, וייבדק מול שאלות האמת & \cite{chen2024m3}}
  \pcrow{אחזור היברידי & חיפוש מילולי \latin{BM25} וחיפוש וקטורי צפוף, ואיחוד הדירוגים בשיטת \latin{RRF}; סינון לפי \texttt{\latin{access\_level}} ו-\texttt{\latin{valid\_to}} & חיפוש מילולי תופס מונחים ומספרי סעיף, והצפוף תופס ניסוח שונה; כשלי אחזור הם מקור עיקרי לכשלי \latin{RAG} & \cite{barnett2024seven}}
  \pcrow{מסד וקטורי ומטא-דאטה & מסד וקטורי מקומי בקוד פתוח (למשל \latin{Qdrant} או \latin{Chroma}) ומסד יחסי מקומי לסכמות פרק~\ref{ch:alt} & כל השירותים כתהליכים מקומיים; האינדקס נבנה מחדש בעת שינוי גרסה & טבלה~\ref{tab:schema}}
  \pcrow{הפקת תצורות ותצוגה & פלט \latin{JSON} לפי \latin{JSON Schema}, ולידציה לפני אישור; תבניות \threeD{} פרמטריות במנוע \latin{WebGL} בדפדפן & אין קומפילציה מחדש בעת שינוי תוכן; בדיקה מוקדמת בדפדפן (\risk{12}) & סעיף~\ref{sec:solgen}}
\end{pctable}
\end{table}
\textbf{אחזור.} הדירוג המאוחד של מסמך \mth{d} הוא סכום ההפכים של הדירוגים בשתי השיטות,
\begin{equation}
  \mathrm{RRF}(d)=\sum_{r\in\{\mathrm{BM25},\,\mathrm{dense}\}}\frac{1}{k+\mathrm{rank}_r(d)},\qquad k=60,
\end{equation}
כאשר \mth{k=60} הוא ערך התחלתי מקובל שיכוונן על שאלות האמת (\latin{D2}). למחולל נשלחים חמשת המקטעים המובילים שעברו סינון הרשאה וגרסה; אם הציון של המקטע
המוביל נמוך מסף \mth{\tau}, שנקבע על עשרת מקרי חוסר-הראיות, המערכת מחזירה “חוסר ראיות” במקום תשובה (\req{05}). אורך ההקשר וגודל המקטעים משפיעים
על האיכות ועל הזיכרון \cite{vladika2025context}, ולכן יוגדרו כפרמטרים ויכווננו בניסוי ולא יקבעו מראש.

\section{“הכלי שמייצר את הכלי”: מ-\latin{RAG} לתצורת תרחיש \threeDadj}\label{sec:solgen}
דרישת הליבה של פרויקט~\latin{5} היא כלי שמייצר את כלי ההדרכה ומעדכן אותו אוטומטית כשהידע משתנה, ולא סימולטור שנבנה ידנית פעם אחת. בניית סצנה
\threeDadj{} מאפס לכל נוהל דורשת מפתח ואנימטור ולא ניתנת להרחבה. לכן ההצעה מפרידה בין שני רבדים. \textbf{הרובד הקבוע} הוא ספריית תבניות
\threeD{} פרמטריות, שהצוות בונה פעם אחת (שלוש תבניות בפיילוט: תחנת בדיקה, סדר פעולות ועץ החלטה); כל תבנית חושפת משבצות: אובייקטים, מצב התחלתי,
צמתי החלטה, אפשרויות, תשובה נכונה, משוב ועוגן ראיה. \textbf{הרובד המתחלף} הוא תצורת \latin{JSON} שמנוע ה-\latin{RAG} מפיק מיחידה מאושרת, וממלא את
המשבצות בזמן ריצה. התוצאה היא שמימוש סצנה חדשה הוא הפקת קובץ ואישורו, ולא כתיבת סצנה (איור~\ref{fig:toolchain}).
\begin{figure}[!htbp]
\centering
\pcFig{toolchain}{%
\pcAutoFit{\begin{pctikz}[x=1cm,y=1cm]
  \tikzset{st/.style={pcnode,text width=2.7cm,font=\scriptsize\sffamily,minimum height=1.0cm}}
  \node[st] (n1) at (0,0) {\tn{1. יחידה מאושרת ומקטעי המקור שלה}};
  \node[st] (n2) at (-3.3,0) {\tn{2. מחולל: מודל מקומי מפיק \latin{JSON} תחת סכמה}};
  \node[st] (n3) at (-6.6,0) {\tn{3. ולידציה: סכמה, עיגון ונאמנות}};
  \node[st,fill=pcWarn!15,draw=pcWarn] (n4) at (-9.9,0) {\tn{4. אישור מומחה בתצוגת הפרשים}};
  \node[st] (n5) at (-9.9,-2.1) {\tn{5. תבנית \latin{3D} פרמטרית טוענת את התצורה}};
  \node[st,fill=pcGood!15,draw=pcGood] (n6) at (-6.6,-2.1) {\tn{6. לומד: משימות, החלטות ויומן אירועים}};
  \node[st,dashed,fill=white] (n7) at (0,-2.1) {\tn{שינוי גרסת מקור: \latin{stale} והפקה מחדש}};
  \foreach \a/\b in {n1/n2,n2/n3,n3/n4,n4/n5,n5/n6}{\draw[pcflow] (\a) -- (\b);}
  \draw[pcflow,dashed] (n7.north) -- (n1.south);
  \draw[pcflow,dashed] (n3.north) to[bend left=40] node[above,font=\scriptsize\sffamily]{\tn{כשל: ניסיון חוזר}} (n2.north);
\end{pctikz}}}
\caption{שרשרת “הכלי שמייצר את הכלי”: יחידה מאושרת הופכת לתצורת \latin{JSON}, עוברת ולידציה ואישור ומאכלסת תבנית \threeDadj{}; שינוי גרסה מפעיל הפקה מחדש}\label{fig:toolchain}
\end{figure}
איור~\ref{fig:json} מציג תצורה מדומה אחת, ואחריה כללי הוולידציה שהמחולל חייב לעבור לפני שהיא מגיעה למומחה.
\begin{figure}[!htbp]
\centering
\begin{english}\begin{minipage}{11.6cm}
\hrule\vspace{3pt}
\fontsize{7.6}{9.2}\selectfont
\begin{verbatim}
{ "scenario_id": "SC-001", "unit_id": "U-014",
  "template_id": "T1_inspection_station", "scene_version": 3,
  "source": { "document_id": "FAA-H-8083-30B", "revision": "rev-B" },
  "scene": { "objects": ["work_card", "revision_tag", "gauge"],
             "initial_state": "revision_unchecked" },
  "nodes": [
    { "node_id": "N1",
      "prompt": "Which document revision applies to this task?",
      "options": ["card_value", "current_revision", "ask_lead"],
      "correct": "current_revision",
      "feedback": { "card_value": "The card may be outdated." },
      "evidence_id": "CH-0412",
      "success_condition": "node_correct" } ] }
\end{verbatim}
\vspace{1pt}\hrule
\end{minipage}\end{english}
\caption{דוגמה מדומה לתצורת תרחיש \latin{JSON} שמאכלסת תבנית \threeD{}; השדות תואמים לישות \texttt{\latin{ScenarioConfig}} בטבלה~\ref{tab:schema}}\label{fig:json}
\end{figure}
\begin{enumerate}
  \item \textbf{תקינות מבנית:} התצורה מתאימה ל-\latin{JSON Schema}: שדות חובה קיימים, \texttt{\latin{template\_id}} מוכר, \texttt{\latin{correct}} הוא אחת האפשרויות.
  \item \textbf{עיגון:} לכל צומת \texttt{\latin{evidence\_id}} שקיים במאגר, ושהמקטע שלו בתוקף (\texttt{\latin{valid\_to}} ריק).
  \item \textbf{נאמנות:} כל משפט משוב נבדק כטענה מול המקטע המצוטט, באותה בדיקה של \req{03}.
  \item \textbf{תנאי הצלחה:} לכל משימה מוגדר \texttt{\latin{success\_condition}}, כדי שהסימולציה תקבע בעצמה אם הושלמה (\req{15}).
\end{enumerate}
כשל בבדיקה גורר ניסיון חוזר אחד עם הודעת השגיאה, ואחריו העברה לעריכה אנושית. אישור המומחה נשאר שער הפרסום.

\textbf{עדכון גרסה אוטומטי.} כששינוי גרסה מתגלה בצינור הקליטה (איור~\ref{fig:pipeline}), המקטעים הישנים נסגרים, הטענות התלויות הופכות ל-\texttt{\latin{stale}}, והיחידות והתצורות
שנשענות עליהן מושהות (\req{11}). המחולל מפיק תצורה חדשה מן המקטעים התקפים, והמומחה מקבל תצוגת הפרשים שמסמנת רק את הצמתים ששונו; לאחר אישור עולה
\texttt{\latin{scene\_version}}, וסשנים שכבר התחילו מסתיימים בגרסה הישנה. האוטומציה היא ביצירה ובהשהיה, והפרסום נשאר באישור אנושי. משך העדכון, משינוי המקור
ועד תצורה מאושרת, יימדד על חמשת שינויי הגרסה המדומים (\latin{D5}) ויושווה לעדכון ידני של סצנה; יעד מספרי ייקבע לאחר מדידת הבסיס.

\section{אבטחת מידע וסביבת פריסה}\label{sec:solsec}
הפיילוט האקדמי ירוץ על מידע פומבי בלבד. במימוש ארגוני עתידי יידרשו סיווג מידע, אישור לייבוא מסמכים, הפרדה בין משתמשים לפי הרשאה, יומן אירועים, מחיקת
מסמכים מהאינדקס לאחר ביטול הרשאה והקפדה על תקני אבטחת מידע רלוונטיים \cite{iso27001}. גם מודלי \latin{3D}, תסריטי הסתעפות והקלטות פעולות ייחשבו נכסי מידע,
ואין לבנות דגם נאמן של ציוד מסווג. האילוצים עצמם מופיעים בטבלה~\ref{tab:constraints}; כך הפתרון עומד בהם:
\begin{itemize}[nosep]
  \item \textbf{סיווג והרשאות:} מידע פומבי בלבד בפיילוט; שדה \texttt{\latin{access\_level}} לכל מקטע, והמנוע אינו ניגש למסמך שאינו מותר לתפקיד.
  \item \textbf{עיבוד מקומי (\req{09}):} אחזור, מסד נתונים, ממשק מומחים ומודל שפה פועלים ללא תקשורת יוצאת; בדיקת רשת בכל ריצה, וניסיון תקשורת חוץ עוצר את הריצה (\risk{02}).
  \item \textbf{דליפה ותיעוד (\req{10}):} יומן מקורות, אישורים וגרסאות, ותיעוד \pct{100} מפעולות העריכה; סיכוני דליפה באחזור נבחנים לפי \cite{zeng2024good}; יחידות תלויות נחסמות עד אישור מחודש.
\end{itemize}
זהו יעד ארכיטקטוני שיחייב בדיקת רשת, ולא הצהרת התאמה פורמלית למדיניות התעשייה האווירית.

\section{מקרה בוחן מייצג}\label{sec:solcase}
נקודת המוצא היא המדריך \latin{FAA-H-8083-30B}, המסביר עקרונות כלליים בהכשרת טכנאי תחזוקת תעופה \cite{faa2023amt}. נבחר תרחיש חינוכי ללא ציוד אמיתי:
טכנאי מתלמד מזהה אי-התאמה בין נתון בכרטיס עבודה לבין מהדורת מסמך העזר. בסצנה גנרית מוצגת עמדת בדיקה עם שלושה כרטיסים: נתוני פעולה מדומים, מזהה גרסת מסמך
ומד מצב (טבלה~\ref{tab:case}). בחירה שגויה משנה רק את מצב הסצנה המדומה ומציגה הסבר, ואין בתרחיש הוראות תחזוקה אמיתיות או ערכי מומנט.
\begin{table}[!htbp]
\centering
\pcCompactTables
\caption{מקרה בוחן מייצג שמורץ לאורך כל השיטה}\label{tab:case}
\begin{pctable}{c,H{0.24},H{0.30},H{0.26}}
  \pchead{שלב & קלט & פעולה & פלט}
  \pcrow{\latin{1} & סצנה עם שלושה כרטיסים & איתור רכיב הבקרה הרלוונטי ושאלה: אילו מסמכים מצביעים על אופן אימות גרסה? & רכיב מזוהה ורשימת מסמכים}
  \pcrow{\latin{2} & שאלת הלומד & פתיחת גרסת המקור המעוגנת במערכת ובחירת המקור התקף & מקור וגרסה עם סימוכין}
  \pcrow{\latin{3} & שתי גרסאות מסמך & השוואת גרסאות, זיהוי הנתונים החסרים ותיעוד המסמך שיש לעדכן & רשימת פערים ורשומת תיעוד}
  \pcrow{\latin{4} & מצב לחץ מדומה המסתיר נתון לא חיוני & הפניית השאלה לגורם מוסמך לפני פעולה פיזית & החלטה, זמן תגובה ומשוב}
\end{pctable}
\end{table}
\textbf{מסלול גיבוי:} אם הגישה העיקרית לא תעבוד, נממש את מסגרת ההערכה של מאמר הדגל \cite{hua2025civask} כפי שהיא (פרק~\ref{ch:lit}): עוזר למידה מבוסס אחזור על הקורפוס הפומבי, עם שאלות בעלות תשובת אמת ובדיקת אחזור, ללא שכבת ה-\latin{3D}.

\section{הנתונים}\label{sec:soldata}
קורפוס הפיילוט המתוכנן יכלול כ-\rng{30}{60} עמודים נבחרים מתוך מדריך פומבי, לפחות \latin{20} שאלות שהכין מעריך אנושי, \latin{10} מקרי חוסר-ראיות,
\latin{5} שינויי גרסה מדומים, \latin{10} שאלות תרחיש וסצנת \latin{3D} גנרית אחת עם שלוש נקודות החלטה. המספרים הם תכנון מדגם ולא תוצאות. כל שאלה תתועד ב“כרטיס נתון” (מזהה, נוסח, מסמך, גרסה, עמוד, תשובת אמת מאושרת, קושי, תוצאת המודל ושם המאשר) ללא פרטי עובד מזהים.
\begin{table}[!htbp]
\centering
\pcCompactTables
\caption{מלאי הנתונים: מה יש, מאיפה ומה חסר}\label{tab:data}
\begin{pctable}{c,H{0.19},H{0.16},C{0.12},H{0.17},H{0.18}}
  \pchead{מזהה & מערך נתונים & מקור ובעלים & סטטוס & היקף ופורמט & מגבלות ושימוש}
  \pcrow{\latin{D1} & מדריך הכשרה פומבי של \latin{FAA} & אתר \latin{FAA} \cite{faa2023amt} & זמין & \rng{30}{60} עמ', \latin{PDF} & קורפוס עיקרי; \idlist{REQ-01, 02}}
  \pcrow{\latin{D2} & שאלות אמת ומקורות מצוטטים & צוות \latin{OrbiTech} & ייווצר & \latin{20} שאלות + \latin{10} חוסר-ראיות & בדיקת אחזור; \hyp{3}, \req{05}}
  \pcrow{\latin{D3} & ראיונות מומחים לא מסווגים & מומחים בהסכמה & טרם נאספו & ראיונות מובנים & ידע סמוי; \hyp{1}, \req{16}; תלוי באישור}
  \pcrow{\latin{D4} & זמני הכנת יחידות ידנית ודיוק חילוץ בסיס & מדידת הצוות & טרם נמדדו & טופסי זמן ושיפוט & בסיס השוואה; \hyp{1}, \hyp{2}}
  \pcrow{\latin{D5} & גרסאות ושאילתות סינתטיות & צוות \latin{OrbiTech} & ייווצר & \latin{5} שינויי גרסה & בדיקת חוסן ועדכון; \req{11}}
  \pcrow{\latin{D6} & תבניות \threeD{} פרמטריות ותצורות \latin{JSON} & צוות \latin{OrbiTech} & תפותח & \latin{3} תבניות, \latin{3} צמתים & המחשת מצב; \req{13}}
  \pcrow{\latin{D7} & יומני בחירה וזמני תגובה & משתתפי הפיילוט & ייאספו בהסכמה & אירועים בסימולציה & כשירות; \hyp{4}, \req{14}}
\end{pctable}
\end{table}
\textbf{מה חסר ואיך נסגור את הפער:} חסרים נתוני בסיס ארגוניים (משך חפיפה, שיעור עזיבה) ומסמכים פנימיים. הבקשה תועבר דרך איש הקשר בתעשייה האווירית (פרק~\ref{ch:summary});
אם לא תאושר גישה לראיונות מומחים עד שבוע \latin{4}, יודגם רק מסלול יצירת הידע ממסמכים פומביים, ומסלול הראיונות יוצג כתכן עתידי שלא אומת (\risk{07}).
\textbf{נתונים מדומים (אם הארגון טרם מסר נתונים).} ההנחות מוצהרות מראש: חמישה שינויי גרסה במסמך, כי שינוי גרסה הוא אירוע שכיח בנהלים, ולכן נבדק ערבוב גרסאות
\cite{zheng2025knowshift}; עשרה מקרים של שאלות ללא מקור, לבדיקת הימנעות נכונה (\req{05}); מצב עומס מדומה (הסתרה זמנית של נתון לא חיוני) לתרגול מבוקר בלבד \cite{nist80084};
וערכים מדומים בלבד בכרטיסי הסצנה, ללא הוראות תחזוקה אמיתיות.
\textbf{קורפוס תחליפי פומבי:} \latin{FAA-H-8083-30B} \cite{faa2023amt} מייצג את מבנה המסמכים הארגוניים (פרקים, טבלאות, איורים והפניות) בלי לחשוף מידע מסווג. בכל תוצר יסומנו נתוני
אמת לעומת נתונים מדומים.

\section{תוכנית הבדיקה והתוצרים}\label{sec:soleval}
טבלה~\ref{tab:eval} קושרת כל דרישה והשערה למדד, לבסיס, ליעד ולקריטריון הצלחה. ערכי הבסיס והיעדים הם הנחות תכן ויעדי תכן, ואינם תוצאות שנמדדו בארגון.
\begin{table}[!htbp]
\centering
\pcCompactTables
\caption{תוכנית הבדיקה: איך נדע שהפתרון עובד}\label{tab:eval}
\begin{pctable}{C{0.14},H{0.22},C{0.15},C{0.15},H{0.34}}
  \pchead{דרישה / השערה & מדד ותרחיש & בסיס & יעד & קריטריון הצלחה}
  \pcrow{\req{16} / \hyp{1} & דיוק חילוץ ידע מראיונות & \pct{60} (הנחה; יימדד) & \gesym{\pct{90}} & שני מעריכים מול מקור תקף; הפרכה אם נמוך מהיעד}
  \pcrow{\req{04} / \hyp{2} & זמן חציוני להפקת יחידה מאושרת & הכנה ידנית \latin{(D4)} & \lesym{\pct{50}} מהבסיס & חיסכון \pct{50} לפחות במשימות שקולות}
  \pcrow{\req{03} / \hyp{3} & נאמנות עובדתית במדגם & --- & \gesym{\pct{90}} & שני מעריכים; אי-הסכמה למומחה שלישי}
  \pcrow{\req{12} / \hyp{3} & איכות שאלות לפי מחוון & --- & \gesym{\pct{80}} & לפחות \latin{6/8} נקודות וללא אפס בנכונות}
  \pcrow{\req{15} / \hyp{4} & שיעור השלמת משימות ללא עזרה בסצנת \latin{3D} & תרגיל \twoD{} זהה & \gesym{\pct{80}} & השלמה: מדד קבלה; הפרש מול \twoD{} \gesym{10} נק' אחוז הוא גישושי}
  \pcrow{\req{05} & הימנעות נכונה בשאלות ללא מקור & --- & \gesym{\pct{90}} & תשובת “חוסר ראיות” במקום תשובה מומצאת}
  \pcrow{\req{08} & זמן שליפת מקור לשאילתה & --- & חציון \lesym{10} \latin{s} & בדיקת קצב בשבוע \latin{5}}
  \pcrow{\req{11} & ביטול פרסום של יחידות תלויות & --- & עד יום עבודה & \latin{5} שינויי גרסה מדומים \latin{(D5)}}
  \pcrow{\req{09} & בקשות לרשת חוץ & --- & \latin{0} & בדיקת רשת נקייה בכל ריצה}
\end{pctable}
\end{table}
\textbf{מסגרת ההערכה:} המדדים מוגדרים מראש ומחושבים כך:
\begin{gather*}
  \hemath{חיסכון בזמן} = 100\cdot\Bigl(1-\frac{T_{RAG}}{T_{manual}}\Bigr),\qquad
  \hemath{נאמנות} = 100\cdot\frac{N_{supported}}{N_{claims}},\\
  \hemath{הימנעות נכונה} = 100\cdot\frac{N_{correct\ abstentions}}{N_{unanswerable}},\qquad
  \hemath{דיוק חילוץ} = 100\cdot\frac{N_{verified}}{N_{extracted}}.
\end{gather*}
בתרחיש \threeD{}: שיעור השלמה \mth{=N_{completed}/N_{started}} ודיוק החלטות \mth{=N_{correct}/N_{decisions}}. ההערכה תתבצע על שאלות שהוגדרו מראש, בידי בוחן אנושי
שאינו שופט את עבודתו שלו; מחוון השאלות כולל ארבעה ממדים בסולם \rng{0}{2}: התאמה למטרה, נכונות תשובה, הסבר ומקור, והוגנות מסיחים. לשחזור הניסוי יישמרו גרסת הקורפוס,
רשימת השאלות, מקטעי האחזור, גרסת המודל והקוונטיזציה, הגדרות החיפוש, גרסת הקוד ויומני הזמן; לכל פעולה בסצנה יירשם אירוע לפי ישות \texttt{\latin{EventLog}}
(טבלה~\ref{tab:schema}), עם מזהה נבדק אקראי בלבד.\par
\textbf{תוצרים:} אב-טיפוס פעיל ראשון (אינדקס ואחזור על הקורפוס הפומבי) בשבוע \latin{5}, ואב-טיפוס \latin{RAG} ו-\latin{3D} לחיץ בשבוע \latin{9}. בכנס הפרויקטים תוצג
הדגמה ניתנת לשחזור שבה משנים מסמך או גרסה ורואים את היחידה, השאלות והתרחיש מתעדכנים. יימסרו גם קורפוס פומבי עם מקור ורישיון, סט שאלות מתועד, יומן בחירות,
גיליון מדדי זמן ונאמנות, מסמך סיכונים ותוכנית הרחבה מותנית אישורים (פרק~\ref{ch:plan}).
\pcChapterEnd{solution}
